%\documentclass[11pt, draft, oneside]{report}
\documentclass[11pt]{report}
\makeindex
\usepackage{ngerman,a4,showidx} %index an seite
%\usepackage{ngerman,a4wide}
%for pdfLaTeX output Bilder als .png speichern:
\usepackage{listings}

\usepackage[pdftex]{graphicx}
% f\"{u}r normales LaTeX->dvi, Bilder als .eps speichern:
\usepackage{graphicx} \DeclareGraphicsExtensions{.eps} %\graphicspath{{bilder/eps/}}

%Definition der Seitengr"osse
\setlength{\textwidth}{15 true cm}
\setlength{\textheight}{22 true cm}
\oddsidemargin  0.5 cm
\evensidemargin 0.5 cm
\topmargin      0 cm

\selectlanguage{german}
%Beispiel fuer ein neues LaTex Kommando
\newcommand{\QSIM}{{\sc QSim}}


\begin{document}

\begin{titlepage}
\begin{center}
  \vspace*{0.5cm}
  {\LARGE Report Studiomesstechnik LU - Group 6} \\
  \vspace{15mm}
  {\huge \bf Laboratory Day 3 - Measurement of DACs and ADCs \\}

  \vspace{15mm}
  {\LARGE Andreas Johann H\"ormer} \\
{\small member 1\\
member 2\\
member 3\\
member 4\\
member 5\\ }
  \vspace{10mm}%15
  -------------------------------------- \\
  \vspace{10mm}%15
  \large
  Institute for signal processing and speech communication \\
  Graz University of Technology \\


  %Vorstand: O.\,Univ.-Prof.\,Dipl.-Ing.\,Dr.\,techn.\,Reinhold Wei{\ss} \\
  \vspace{15mm}%1
  \begin{figure}[!ht]
  \begin{center}
  \centerline{\includegraphics[width=4cm,keepaspectratio=true]{TULogoneu}}
  \end{center}
  \end{figure}
  \vspace{10mm}
Supervisor: DI Wolfgang Truppe \\
  %Begutachter: O.\,A.o.Univ.-Prof.\,Dipl.-Ing.\,Dr.\,techn.\,Eugen Brenner \\
  %\vspace{5mm}
  %Betreuer: Univ.-Ass. Dipl.-Ing.\,Dr.\,techn.\,N. N.\\
  \vfill
  %\newline
  %\normalsize
  audio studio of TU Graz, 11.03.2015
  \vspace{0.5cm}
\end{center}
\end{titlepage}

% Die Titelseite ist immer in Deutsch (austrian), danach h\"{a}ngt es von der
% Sprache der Diplomarbeit ab. Jedenfalls muss eine Kurzfassung und
% ein Abstract existieren

%\thispagestyle{empty} 
%\selectlanguage{english}


\newpage
\selectlanguage{english}
\vspace*{2.2 cm}
{\Large
\noindent
{\bf Abstract}} \\
\vspace*{0.3 cm}

\noindent
Main target of this laboratory was the measurement of different DACs and ADCs of the fully digital mixing panel LAWO $mc^2 66$. Additionally preamplifiers of a high quality input compared to standard inputs were measured. Measurements were done qualitatively in terms of dynamic ranges and frequency characteristics.\\
This report is consisting of XXXXXXXXXXXXXX pages. 


%\newpage


%\selectlanguage{english}
%\vspace*{2.2 cm}
%{\Large
%\noindent
%{\bf STATUTORY DECLARATION}} \\
%\vspace*{0.3 cm}

%\noindent
%I declare that I have authored this thesis independently, that I have not used other than the declared sources / %resources, and that I have explicitly marked all material which has been quoted either literally or by content from %the used sources.
%\vspace*{0.3 cm}

%\vspace{2 cm}

%\noindent ..............................\hfill ...........................................


%\noindent date  \hfill (signature)


\newpage
\selectlanguage{english}
%\newpage
%\vspace*{2.2 cm}
%{\Large
%\noindent
%{\bf Danksagung}} \\
%\vspace*{0.3 cm}
% OPTIONAL

%\noindent
%Diese Diplomarbeit wurde im (Studien)Jahr am Institut f\"{u}r
%Technische Informatik an der Technischen Universit\"{a}t Graz
%durchgef\"{u}hrt.

%\smallskip
%Danksagung an alle am Institut bzw. bei Firmen, die geholfen
%haben....

%\medskip
%Danksagung an Freunde und Freundinnen f\"{u}r das Verst\"{a}ndnis, ebenso
%den Eltern und allen sonstigen Sponsoren....

%\vspace{2 cm}

%
%\noindent Graz, im Monat Jahr \hfill Name des Diplomanden

\newpage
% Inhaltsverzeichnis
\tableofcontents  

% Tabellenverzeichnis
% OPTIONAL
%




\listoffigures 

% Abbildungsverzeichnis
% OPTIONAL
%\listoftables

%Seitennummerierung am Kopf inkl. Kapitel"uberschrift
\pagestyle{headings}

%---------------------------------------------------------------------------------------------------------
%---------------------------------------------------------------------------------------------------------
\chapter{Preamble}
% Einleitung ist Kapitel 1
\label{kap:Einleitung} 

%---------------------------------------------------------------------------------------------------------
\section{Motivation}

%---------------------------------------------------------------------------------------------------------
\section{Purpose}



%---------------------------------------------------------------------------------------------------------
%---------------------------------------------------------------------------------------------------------
\chapter{DAC measurements}
\section{Introduction}
\section{Measurements}
	\subsection{Experimental setup}
		\subsubsection{Equipment}
	\subsection{Results}
	\subsection{Discussion}


%---------------------------------------------------------------------------------------------------------
%---------------------------------------------------------------------------------------------------------
\chapter{ADC/preamplifier measurements}
\section{Introduction}
\section{Measurements}
	\subsection{Experimental setup}
		\subsubsection{Equipment}
	\subsection{Results}
	\subsection{Discussion}

%---------------------------------------------------------------------------------------------------------
%---------------------------------------------------------------------------------------------------------
\chapter{Conclusion}



\begin{center}
\begin{tabular}{|c||c|c|}
\hline 
SIL & 	PFH &		RRF\\ \hline
1 &	0.00001-0.000001 &	100,000-1,000,000\\
2 &	0.000001-0.0000001 &		1,000,000-10,000,000\\
3 &	0.0000001-0.00000001 		& 10,000,000-100,000,000\\
4 &	0.00000001-0.000000001 &	100,000,000-1,000,000,000\\
\hline
\end{tabular}
\end{center}


\begin{center}
\begin{figure}[htbp]
\includegraphics[width=11cm,keepaspectratio=true]{PROFIsafe}
\caption{PROFIsafe}
\label{Abb.:1}
\end{figure}
\end{center}


\begin{appendix}
\chapter{List of Abbreviations}
\chapter{High resolution diagrams}
\end{appendix}

%nicht referenzierte Literaturstellen

\nocite{mseifter88,pmandl97,weiss92,weiss92a,ginthoer93}

% Literaturverzeichnis einbinden, alpha, plain, unsrt, abbrv
\newpage
%Eintrag im Inhaltsverzeichnis
\addtocounter{page}{1}
\addcontentsline{toc}{chapter}{Literaturverzeichnis}
\addtocounter{page}{-1}

\bibliographystyle{alpha}

\bibliography{protokoll}

%Index File
%\input{diplom.ind}

\end{document}
